\section{Method: the \sys\ Pipeline}
\label{sec:method}

\subsection{Design principles}
\label{sec:principles}

\begin{enumerate}
  \item \textbf{Focus on computable subfields.} Prioritize computational biology
        and bioinformatics, where papers describe algorithms, models, or
        analysis pipelines with standardized inputs and outputs.
  \item \textbf{Automation-first, expert-review-as-backstop.} The multi-agent
        pipeline generates environments; human experts review a subset for
        fidelity and shortcut resistance.
  \item \textbf{Standardized output.} All environments are packaged as \harbor\
        tasks supporting containerized execution, automated scoring, and result
        storage.
  \item \textbf{Verification before packaging.} No environment is published
        unless it passes the four gates of \S\ref{sec:gates}. A missing
        dependency or an absent grader reports \textsc{Unverified}, never a pass.
\end{enumerate}

\subsection{Pipeline architecture}
\label{sec:architecture}

\begin{quote}
Literature Acquisition $\rightarrow$ Paper Parsing $\rightarrow$ Environment
Construction $\rightarrow$ Baseline Reproduction $\rightarrow$ Quality Review
$\rightarrow$ \harbor\ Packaging $\rightarrow$ Evaluation \& Iteration
\end{quote}

\paragraph{Curator Agent (new).}
Screens Europe PMC, PMC, and bioRxiv for papers meeting computability criteria:
open-access full text, computational biology/bioinformatics classification,
available code or data, and defined evaluation metrics (F1, AUROC, RMSE, etc.).
Target: 1,000 papers screened, 200 selected. In the released framework the
screening rules are deterministic (licence, publication year, abstract length,
subfield keyword); a model is consulted \emph{only} for boundary papers that
pass every rule except the computational classification, so the selection is
reproducible and auditable.

\paragraph{Construction Agent.}
Generates the Dockerfile, dependency installation, task description, input data,
and evaluation script, producing an initial executable environment. Its output is
schema-checked against a strict blueprint schema before anything is written to
disk.

\paragraph{Reproduction Agent.}
Independently reimplements the paper's method to establish a baseline score
through a \emph{second channel} that does not see the construction agent's own
claimed numbers; agreement between the two channels---not either channel
alone---is the signal. If the environment fails to run, an automated debugging
loop attempts resolution; persistent failures are flagged for expert
intervention.

\paragraph{Review Agent.}
Checks fidelity to the source paper, resistance to shortcut strategies, baseline
stability, and safety. Biosafety screening is \emph{deterministic and
fail-closed}: a dual-use denylist is applied to the task objective before any
model is consulted, and a match blocks packaging outright (claim~\cid{C3}). An
ambiguous or erroring review is recorded as a finding, never as a silent pass.

\paragraph{Packager Agent (new).}
Converts validated environments into \harbor\ task directories, re-verifying at
the last gate and writing a machine-readable manifest line per environment---the
artifact that Stage~1 tooling consumes.

\paragraph{Human Expert.}
Reviews a curated subset of environments, focusing on methodological fidelity
and evaluator validity.

\subsection{\harbor\ task format}
\label{sec:format}

We target \harbor~0.24's on-disk layout, which the released framework
cross-validates against directly (claim~\cid{C2}):

\begin{verbatim}
tasks/<slug>/
  task.toml               # Harbor TaskConfig, schema_version = "1.4"
  instruction.md          # the task as the agent sees it
  environment/Dockerfile  # pre-configured bioinformatics toolchain
  tests/test.sh           # agent execution entrypoint -> grader
  tests/grade.py          # automated scorer (F1 / AUROC / RMSE / ...)
  solution/               # reference solution
  rubric.md               # expert-authored rubric (optional)
\end{verbatim}

\footnotetext{Earlier descriptions of \harbor\ task sets used a
\texttt{task.yaml} + \texttt{inputs/} + \texttt{run.sh} + \texttt{evaluator/}
layout. We follow the format that \harbor~0.24 actually validates, and our
emitter is checked against the upstream \texttt{TaskConfig} model rather than
against that earlier description.}

\paragraph{Why the format matters for verification.}
\harbor's \texttt{TaskConfig} model ignores unknown keys by default, so a
mistyped resource key (e.g.\ \texttt{cpu} where \texttt{cpus} is meant) is
\emph{silently dropped} and the task runs with unintended limits. Our first
verification gate is therefore an embedded JSON Schema with
\texttt{additionalProperties: false} at every level, which rejects such a key
outright (claim~\cid{C1}).

\subsection{Subfield prioritization}

Initial rollout focuses on subfields with high standardization and clear
evaluation metrics: single-cell RNA-seq analysis; variant annotation and
pathogenicity prediction; protein structure prediction; phylogenetic analysis;
biological image segmentation; and gene regulatory network inference.

\subsection{Safety}

All generated tasks are restricted to \emph{computational analysis}. Dual-use
objectives (gain-of-function, pathogen enhancement, and related terms) are
rejected by a deterministic denylist applied before any model call, and any hit
blocks packaging (claim~\cid{C3}). This is a floor, not a ceiling: it catches
the objective stated in the task text, and expert review remains the backstop.
